\chapter{Literature Review}
\label{ch:literature}

\section{Introduction}

This chapter discusses the current body of literature regarding the construction of a real-time decision support dashboard for prioritization of high-risk flights within airline inventory management. The chapter begins with a presentation of key themes in the discussion, which cover three major areas and can together provide an answer to the problem statement. Namely, these are revenue management and inventory control, business intelligence dashboards, and operational aviation dashboards and real-time analytics. In each of these themes the body of current literature is discussed, and the research gap is identified.

\section{Airline Revenue Management and Inventory Control}
\label{sec:ch2:rm}

In operations research studies, one of the main focuses is airline revenue management (RM), since RM has been proven to impact airlines' outcomes financially in a positive way. Thus, the necessity to conduct an in-depth study of RM emerges. Information related to the issues raised in the context of operations research studies in the area of air transport can be found in \citet{BarnhartBelobabaOdoni2003}.

The topic chosen for this thesis has already attracted considerable attention in academia over time, and there is already enough literature available on the same topic; however, there are open directions that this thesis builds upon. According to \citet{Shaw2016}, the introduction of deregulation in airlines in the late 1970s was instrumental in changing the focus of airline strategies towards price differentiation. In order to achieve differentiation, airline firms needed to implement a system that could help them allocate capacity across multiple fare classes according to their respective demand levels, which was achieved via inventory management controls that were enabled by the relaxation of fare regulations that existed prior to deregulation. The main contribution of Shaw lies in laying down the fundamental reason as to why the field of revenue management science emerged, and how it became essential for any airline firm operating after deregulation.

The seminal work by \citet{McGillVanRyzin1999} in \textit{Transportation Science} offers an overview of developments in the revenue management discipline. The two authors conduct a literature review on quantity-oriented and price-oriented revenue management strategies, describe the history of development of revenue management from its origins as yield management in the aviation industry through to its use in other service sectors, and highlight the remaining research issues related to network revenue management, choice models, and system pricing integration.

\citet{BodilyWeatherford1995} developed a framework for revenue management related to perishable assets, taking into consideration generic and multiple prices yield management as well as customer diversion. It was published in the journal \textit{Omega} and is concerned with the trade-off between encouraging price-sensitive customers' demand by offering them discounts and protecting against any diversion that may occur due to non-price-sensitive customers moving towards lower-priced fare classes.

\citet{LeeHersh1993} present a dynamic programming approach to model airline seat inventory management for more than one seat booking per request. This paper was published in the journal \textit{Transportation Science}. The concept of the balancing act of accepting requests at lower fare but taking care of future high-fare demands through the dynamic programming approach presented by the authors inspires including ``days till departure'' into the composite risk score.

The work done by \citet{Belobaba1987} during his Ph.D. studies at the Massachusetts Institute of Technology on seat inventory management as a stochastic control problem paved the way for what is known as the EMSR heuristic for allocating capacity amongst different fare classes. The problem considered by Belobaba was concerned with the trade-off faced by any inventory controller who needs to accept or reject booking requests from customers at different fare classes depending on the cost-benefit analysis of accepting the bookings versus rejecting them. The probabilistic solution suggested in this study is still highly relevant today since the core logic behind it remains the same and forms the backbone of most revenue management solutions implemented today.

This dissertation work was eventually published by \citet{Belobaba1989} in the \textit{Operations Research} journal under the title of EMSR model of airline seat inventory control. In this paper, Belobaba provided an explicit formulation of a decision problem faced by inventory controllers and solved for a threshold condition under which accepting lower-fare class bookings should become worthwhile. This solution required estimation of demand parameter distributions that could easily be estimated and did not require extensive computations that the optimal control problem would otherwise require. The relevance of this study is that EMSR has remained a widely used methodology in revenue management science until now.

\citet{Kimes1989} was responsible for introducing the more generic term \emph{yield management} to describe the strategy that enables a firm to create value out of limited capacity resources. He identified five preconditions that must hold good if yield management was going to be effective in generating revenues. These included fixed capacity availability, market segmentation capabilities, perishable capacity resources, advance sales possibility, and demand variability. Kimes proved that these prerequisites exist across many industries other than airlines, which were initially the only focus of revenue management studies. Therefore, this paper is very significant in that it paved the way for establishing yield management as a science. Furthermore, this paper is particularly relevant because of the importance that it attributes to inventory management analysts.

\citet{SmithLeimkuhlerDarrow1992} document the experience of implementing and operating yield management at American Airlines and quantify its financial benefits in terms of additional revenues earned through overbooking control, discount seats allocation, and traffic forecasting. This paper is significant due to the following two reasons. Firstly, Smith et al.\ demonstrate clearly how yield management works in practice through various examples and explanations regarding how the system operated successfully to generate extra revenues for the airlines. The authors explain why the system worked well, describing various operational aspects. Secondly, it validates the idea that the combination of analyst oversight and automation leads to effective yield management practice; this aspect lies at the core of the research problem addressed in this thesis.

\citet{Chatwin1999} builds upon previous discrete-time models for overbooking decisions under time-dependent fares and refund policies. He develops a stochastic continuous-time model of overbooking and obtains a threshold policy whereby any booking should be accepted as long as the number of bookings falls below certain thresholds. This policy depends on time to departure and hence captures a realistic aspect of the actual decision process that involves tightening overbooking limits as departure time gets closer. The significance of this paper for this thesis lies in the fact that overbooking risk varies over time, and this motivates incorporating time into the composite risk score.

\citet{SubramanianStidhamLautenbacher1999} extend the scope of yield management modelling by building a joint stochastic decision process for all four critical elements involved in airline operations: seat inventory control, overbooking, cancellations and no-shows. Previous work considered separate models for each component independently; however, this is unrealistic since, in practice, the decisions related to overbooking depend on cancellation and no-show probabilities and vice versa. As a result, this formulation yields a combined decision rule for managing inventory in terms of monotonicity in time to departure and booking level, thus providing an important academic precedent for this thesis, where multi-factor risk is considered in a single composite score.

\citet{PakPiersma2002} give an exhaustive overview of operations research techniques applied in airline revenue management and summarise the existing research contributions for the three major categories of techniques employed in this problem domain: EMSR heuristics, Markovian decision models, and network-based bid-price methods. Although there is plenty of literature devoted to optimising the single-resource airline revenue management problem, only few attempts have been made to extend this technique into multi-resource or network revenue management problems, leaving a vast amount of research potential open to future explorers. The relevance of this paper lies in the fact that it establishes revenue management optimisation as the field of interest for this thesis, highlighting the lack of decision support techniques.

\citet{BoydBilegan2003} examine how the emergence of e-commerce influenced airline revenue management practices in relation to online distribution systems. In their view, revenue management shifted away from a planning exercise to a more tactical activity characterised by high responsiveness to dynamic demand conditions. Furthermore, it became increasingly important to have decision support systems that operate in real-time instead of working on pre-stored snapshot data as before. This paper is cited here due to its relevance in terms of describing how the real-world problem has evolved over the years, as this thesis considers online inventory decisions.

\citet{BarnhartBelobabaOdoni2003} gave an extensive overview of the state-of-the-art operations research techniques applicable to solving various kinds of operational problems associated with air transportation services, including fleet scheduling, crew rostering, airport operations and revenue management. In particular, revenue management was pointed out by the authors to be one of the most mature operational areas of air transportation, with empirical revenue uplift estimated to lie in the range 4--6\%. However, the authors noted that even in this case, there seems to be a discrepancy between model complexity and its implementation in real-world practice, a problem that this thesis seeks to solve.

\citet{WeatherfordKimes2003} evaluate several different forecasting methodologies used in hotel industry revenue management systems. They compare exponential smoothing methods with linear regressions and pick-up techniques on simulated data sets and conclude that none of the methodologies can claim superiority under all conditions of demand variability and forecasting horizon. In addition, simple pick-up methods show comparable forecast accuracy as long as high-quality historical booking data is used as input into the forecasting algorithm. This study is significant to this thesis as many of the recommendations given in that study on forecasting accuracy are employed, particularly in relation to the snapshot-based design of the proposed decision support technique.

\citet{TalluriVanRyzin2004} produce the bible of revenue management, which unites decades-long efforts by various researchers in airlines, hotels, car rental, retail stores and entertainment sectors to develop an overarching theoretical foundation of this science. This textbook includes topics on single and network revenue management, auctions, pricing and overbooking problems together with economic justifications for the practice of differential pricing. For the case of airlines, Talluri and van Ryzin present a consolidation of EMSR heuristics and bid-price control into a unified theoretical framework for inventory control. The textbook is important to this research as the work is grounded precisely in these concepts.

\citet{FiigIslerHopperstadBelobaba2010} address an emerging problem in airline revenue management concerning the shift of fare structures from purely independent demand structures suitable for traditional revenue management methods to a new generation of hybrid systems. The problem is that hybrid fare classes allow the application of differential pricing rules while simultaneously retaining unrestricted bookings, thereby violating a crucial assumption on which traditional RM algorithms relied, namely the separation of demand segments. Therefore, a new optimisation approach is presented in this paper that takes into account this problem explicitly through the buy-down effects model. This paper serves as an example of how commercial realities need to be reflected in models.

Despite the abundance of literature available for RM, this whole body of knowledge focuses mainly on optimisation aspects and leaves aside questions pertaining to analyst interaction with model outputs or design considerations for RM dashboards.

\section{Dashboards, Decision Support, and Business Intelligence}
\label{sec:ch2:bi}

A notable article on real-time business intelligence within the aviation industry is \citet{WatsonWixomHoffer2006}, which talks about the implementation of real-time business intelligence within Continental Airlines. The paper talked about how Continental Airlines shifted from overnight batch reporting to a real-time data warehouse. Moreover, Continental reported tremendous improvement within their revenue, cost and service quality. There are three things identified by the researchers that enabled Continental Airlines to have such successful BI, which include: having a sponsor for the project, investments in the integration of data, and presenting relevant analysis to users. This paper is highly relevant to the present research because the central purpose of this thesis is to develop a tool that enables operations managers in the aviation industry to make proper use of their real-time BI.

Another very important research is \citet{WatsonWixom2007} who talked about business intelligence practices in academics and industry and was published within \textit{IEEE Computer}. There are many interesting points in this paper. Researchers described three stages of business intelligence maturity that include: reporting phase, analytical phase, and advanced analytic phase with decision support. The research also describes organisational factors differentiating successful BI from those that are unsuccessful. Researchers emphasised that technology is not enough and organisations need to ensure that their systems deliver benefits to business through the right organisational processes and tools. The reason why the paper is important is that the design philosophy of this thesis is not to create just a BI tool that would sit on top of the IT infrastructure but will be used by the analysts for delivering business benefit.

\citet{ChenChiangStorey2012} offered an excellent overview of business intelligence and analytics within the last two decades in their paper \emph{Business Intelligence and Analytics: From Big Data to Big Impact}. Researchers discussed the evolution of the discipline and divided it into three generations that include: analytics on structured data, analytics on Web and Social Media data, and analytics on Mobile and Sensor data. Researchers explained how different methodologies and techniques were applied for each generation and concluded that Business Intelligence and Analytics needs to be seen as an ongoing capability rather than just a set of technologies that would solve problems immediately after implementation. It is a very important reading, since the paper allows one to understand what technologies are suitable for building the right BI capability in the airline industry.

A systematic literature review done by \citet{VazquezIngelmo2019} on \emph{Information Dashboards and Tailoring Capabilities: A Systematic Literature Review}, published in \textit{IEEE Access}, addresses the issue of how information dashboards handle the problem of user heterogeneity by having tailoring capabilities. The finding from the reviewed studies is that there is a constant struggle between the ``one size fits all'' way that operational dashboards adopt and the requirements of roles. This research is helpful in developing the prioritization dashboard for this dissertation that will be tailored to one role of operator (Inventory Analyst) with different perspectives of the same information.

According to \citet{Pauwels2009} who have written about dashboards for marketing and services in the \textit{Journal of Service Research}, there are frameworks for looking into the three things involved in dashboards that include data exposed, decisions made, and the process of organization. According to the authors, dashboards make sense only when they influence the decisions made based on the information but not just observing information only. Therefore, the recommendation-oriented Overview tab in the created dashboard ensures that each flight that poses higher risk comes with a recommended decision.

\citet{WixomWatson2001} of \textit{MIS Quarterly} have published an empirical analysis of what determines the success of data warehouse development initiatives. Management support, user involvement, and data quality are listed among implementation issues and are shown to be the driving force of data warehousing. Their results apply also to dashboards: the success of this particular prioritization dashboard will not just depend on the effectiveness of the decision tree models but also on data quality and their integration into existing processes.

This thesis also draws on \citet{YigitbasiogluVelcu2012}, a literature review focused on dashboards in performance management. Literature reviews allow synthesising all the main results that focused on dashboards in performance management. The research was published in the \textit{International Journal of Accounting Information Systems}. Literature reviews allow synthesising all the main results to draw conclusions. In this case researchers concluded that there are many elements making dashboards effective, including data visualisation, user interaction, drill-down and alignment of the dashboard to organisation performance goals. However, there is still a big gap between the theoretical guidelines in building dashboards that people talk about and the reality of how those are implemented in practice. Dashboards tend to become overly complex and include excessive amounts of KPIs that distract user attention rather than helping them to accomplish their task effectively. This paper informs the design of the dashboard developed in this thesis, clarifying what should and should not be included.

\citet{Few2013} provides a practitioner reference on dashboard design, defining a dashboard as a single-screen visual interface showing the most important information required to achieve certain objectives. The principles described (avoiding visual noise, conveying magnitude through position, length, and colour, and favouring condensed summary over multi-step reports) informed the design choices made in the dashboard developed for this thesis.

The empirical study by \citet{Rouhani2016}, published in the \textit{Journal of Enterprise Information Management}, focused on evaluating the effect of BI system usage on decision support and benefits gained through it. The authors found out that only those BI systems offering timely, contextually appropriate information produce better decisions. This conclusion is very important because this is one of the main ideas that needs to be considered during designing the dashboard for the operations manager; in particular, it is necessary to avoid general performance measures and focus on creating actionable suggestions.

In addition, \citet{Sarikaya2019} readdressed the question of what is meant by the term ``dashboards''. Their paper entitled \emph{What Do We Talk About When We Talk About Dashboards?} was presented at \textit{IEEE Transactions on Visualization and Computer Graphics}. According to the authors, terminology around dashboards became ambiguous, leading to neglect of important principles behind them. In order to tackle this challenge the researchers suggest grouping dashboards along four axes including: type of the user, type of information provided, level of interactivity, and role in the analytics process. All these characteristics are extremely helpful in evaluating the effectiveness of the developed dashboard and justifying its design decisions.

\section{Aviation Operational Dashboards and Real-Time Analytics}
\label{sec:ch2:aviation}

Based on the findings of general studies in the area of dashboards and BI, a more targeted stream of research has examined how operational dashboards and real-time analytics have been used specifically in the field of aviation. The relevance of such a study to the current project is defined by the fact that it describes the peculiarities of operations associated with fast changing information, quick decision making, and urgent actions based on the received information.

After reviewing the literature on dashboards and business intelligence in general, the review turns to literature specifically related to aviation to try to understand why no such solution is yet available in the market. Such a dashboard becomes particularly important when political developments affect airlines, as aviation operations are full of situations where data changes rapidly and decisions must be made equally quickly. The following section reviews the studies analysed during this research.

\citet{RebolloBalakrishnan2014} present a case in which they employ Random Forest classification to analyse and predict the delays associated with air traffic delays. It has been found out from their study that Random Forest classifiers can recognise delay patterns which cannot be recognised by conventional statistics. This research article has been published in \textit{Transportation Research Part C}. Random Forest classifiers have also been employed in forecasting of this thesis.

The article of \citet{RakshitKrishnamurthyYu1996}, published in the journal \textit{Interfaces}, discusses a tool called System Operations Advisor (SOA) which United Airlines started deploying in 1993 and continued operating until 1994 in its Operation Control Centre to deal with flight disruptions. This application suggested several recovery strategies including delaying, aircraft swapping, and rerouting flights based on real-time operations data collected by the airline. It was estimated that between October 1993 and March 1994 this application saved 27,000 minutes of delays and helped to avoid some costs. The relevance of this article lies in the fact that the study shows the possibility of applying decision-making systems and achieving quantifiable benefits for the airline, as well as establishing a tradition of using a human-in-the-loop approach when it comes to making a final decision.

\citet{Fiig2016} analyse the concept of dynamic pricing to be considered as the next step of the future generations of revenue management systems, noting that traditional RMS use historical data as a basis of their decisions and do not take into account the competitive environment directly. The authors state that the future revenue management should consider dynamic, personalised and competitive pricing policies. Although the priority dashboard designed for this thesis works on inventory decisions and not on pricing decisions, the operational difficulties mentioned by Fiig et al.\ are worth mentioning.

\citet{ClausenLarsenLarsenRezanova2010} in their article which was published in \textit{Computers and Operations Research} discussed the concept of disruption management in the airline industry by reviewing the operations research literature. In the beginning, the authors distinguish between two processes involved in managing unexpected problems in the airline business, scheduling and disruption management, which occur in different time periods and require distinct decision support solutions: disruptive situations usually happen hours or days prior to departure while the scheduling process happens months ahead. In their work the authors provide evidence that is important to this thesis in demonstrating that disruption management is the type of operation that requires real-time decision-making and cannot be considered a scheduling process.

The paper presented by \citet{Pedregosa2011}, published in the \textit{Journal of Machine Learning Research}, introduces the popular machine learning library scikit-learn and describes the main aspects of its implementation such as consistent estimator interface and integration into the scientific Python stack. This article provides the foundation for understanding the basic concepts behind implementing the machine learning models used in the predictive engine of this thesis and in testing the functionality of the tool on the synthetic data.

\citet{Zanzotto2019} published a viewpoint article entitled \emph{Human-in-the-loop Artificial Intelligence} in the \textit{Journal of Artificial Intelligence Research}. The paper criticises current approaches to developing machine learning models by challenging the assumption that such systems could completely replace human decision-making capabilities with AI technology and argues for a different paradigm called human-in-the-loop. According to Zanzotto, the human-in-the-loop approach has at least three key advantages, namely maintaining responsibility for results on the side of the user, integrating tacit knowledge into the decision-making process, and ensuring transparency. As far as the present work is concerned, this article is crucial because the same methodology is followed in the project design.

\citet{MosqueiraRey2023} publish their review article in \textit{Artificial Intelligence Review} discussing the progress in applying human-in-the-loop approaches in the field of machine learning. The review covers several categories of HITL applications according to the stages where human interaction occurs, data labelling, training model, testing, deployment and usage, as well as types of interactions including active learning, machine teaching, explainable AI and interactive ML. This paper is important as it provides valuable background regarding current trends and approaches in the area and helps identify the approach to adopting HITL methodology.

\citet{Salleh2023} discuss the development of an aviation performance dashboard in order to represent information regarding scheduling and reasons for delays in 2020 due to the ongoing COVID-19 pandemic. The authors suggest a methodology that consists of four major phases of development: data collection and cleaning, exploratory analysis and visualisation phase. They emphasise that aviation dashboards are supposed to facilitate operational and tactical decision-making of managers involved in monitoring operations. This paper is relevant to the present research as the authors propose a practical approach to designing such dashboards which aligns with the methodology adopted in this thesis.

To conclude, the review of existing literature on operational decision-making and real-time analysis proved the viability of such an application for aviation business operations and its benefits. However, most studies focused mainly on certain problems related to aircraft operations such as scheduling and delay analyses rather than developing general frameworks that can be applied in different cases. This thesis addresses this gap in applied operational management science.

\section{Research Gap}
\label{sec:ch2:gap}

The reviewed literature in all three previous themes that include airline revenue management, dashboards and business intelligence, and operational analytics in aviation reveals one clear trend that needs to be highlighted. The literature for every theme is well represented and extensive enough, yet it seems that little to no research was done on how those areas can overlap. Especially relevant is how the use of dashboards could assist the prioritising process in airline revenue management.

The theme of airline revenue management has gained much attention over decades as multiple mathematical models were built and presented by different authors. Some of the most famous papers published include work of \citet{Belobaba1987, Belobaba1989}, \citet{LeeHersh1993}, \citet{PakPiersma2002}, and more recent contributions such as \citet{FiigIslerHopperstadBelobaba2010}. It appears that every article is aimed at providing better mathematical and computational modelling solutions in order to make decision-making more precise for the airline revenue manager. Yet there were hardly any articles found which would talk about the problems faced by analysts during this work and provide solutions on how to deal with such problems as prioritising and limited amount of time for decision-making for multiple flights.

Some studies have already been made related to dashboards and business intelligence, where authors state that good dashboard development and implementation may increase the efficiency and productivity of decisions. Some even proved through their research that application of real-time business intelligence will bring numerous benefits financially and managerially to the organisation \citep{Andronie2015, YigitbasiogluVelcu2012, WatsonWixomHoffer2006, Rouhani2016}. Nonetheless, almost all research papers are general enough regarding dashboards usage, organisational benefits or even whole strategies on enterprise-wide BI implementation. Only few papers address the issue of using dashboard technology for certain purposes related directly to business intelligence decision making. None of the research was performed considering how the use of dashboards might help solve the prioritising problem that exists in the airline industry.

There exist several papers in scientific literature dedicated to the question of operational decision making in the aviation field by means of implementing dashboard technology. For example, papers published by \citet{RakshitKrishnamurthyYu1996} and more recently by \citet{Salleh2023} show high potential for this technology usage in aviation operations. However, no relevant literature was found that addresses the prioritisation problem in airline revenue management.

Despite this being an important issue, the question requires further research. That is why the aim of this thesis is to create a real-time dashboard tool to assist the inventory analysts in prioritising the flights requiring the most urgent actions prior to departure.

\section{Key Takeaways}

The key findings from the literature reviewed in this chapter are summarised below:

\begin{itemize}
    \item Revenue management has emerged as one of the most studied topics in airline literature. Moreover, applying mathematical and computer-based models for seat inventory control, overbooking policies, pricing decisions, and demand forecasting has proved to positively influence airline revenues by as much as 4--6\% according to studies conducted in \citet{Belobaba1987, Belobaba1989}, \citet{PakPiersma2002}, and \citet{BarnhartBelobabaOdoni2003}.
    \item The decisions made through the process of revenue management analysis are highly time-sensitive, since there comes an hour for decision-making in which any changes become impossible to implement. Therefore, analysts are expected to work under pressure to prioritise flights to ensure timely completion of their jobs.
    \item Dashboards and business intelligence applications are widely recognised as tools facilitating decision-making processes in organisations. More importantly, real-time business intelligence has shown to deliver significant financial and operational gains for airlines, proving its worth with an impressive ROI of 500\% achieved through investing only 30 million dollars into this technology at Continental Airlines \citep{WatsonWixomHoffer2006}.
    \item Real-time decision making supported by operational aviation dashboards has proven valuable across many areas including aircraft operations controls, air traffic management, delay analysis and others. One of those examples is System Operations Advisor introduced by United Airlines that saved the company nearly half a million dollars in costs incurred because of flight delays associated with the operations control process \citep{RakshitKrishnamurthyYu1996}. The same application managed to reduce delays in this particular domain by 50\%.
    \item Although each stream is advanced, nothing prior brings the two streams together in an integrated tool to assist revenue management in making a decision, and that is what this thesis aims to achieve.
\end{itemize}
